A general framework has been developed for formulating cross-sectional constitutive models that relate the family of objective beam theories in \cref{sec:genwarp} to a three-dimensional continuum for any number of auxiliary warping fields \(\bm{\alpha}\), including none at all. 

In \cref{sec:section-kinematics} three compatibility relations have been derived relating the deformations \(\bm{e}\) to a material strain tensor \(\StrainTensor\): the exact form \eqref{eq:exact-gl-strain-Phi}, the partial truncation \eqref{eq:wagner-gl-strain-pi}, and the total linearization \eqref{eq:linear-gl-strain-pi}. 
All three representations of the material strains may be employed in connection with a finite deformation beam theory \eqref{eq:finite-alignment-field}, or the linearization \eqref{eq:trace-lift-linear}. 
In the remainder of this work attention will be limited to the moderate-strain Wagner model \eqref{eq:wagner-gl-strain-pi}, and the linearized model \eqref{eq:linear-gl-strain-pi}.
These results specialize to several notable works in the literature:
\begin{itemize}
    \item \cite{simo1991geometricallyexact} employed the linearized material kinematics \eqref{eq:linear-gl-strain-pi} in connection with the geometrically exact \eqref{eq:finite-alignment-field} with one warping mode associated to Saint-Venant torsion function. \cite{gruttmann2000theory} generalized \cite{simo1991geometricallyexact} to consider sections in arbitrary coordinate systems and elasto-plastic material response. 
    The present treatment supports an arbitrary number of warping shapes, and allows for accurate incorporation of the Wagner effect in the geometrically exact setting \eqref{eq:finite-alignment-field} by switching the cross sectional model. 
    \item The higher-order strain field \eqref{eq:wagner-gl-strain-pi} was considered by \cite{battini2002corotational-thesis} and \cite{lecorvec2012nonlinear} with the linearized beam theory \eqref{eq:trace-lift-linear}. 
    However, these works employed a constrained continuum derivation which tightly coupled the effect of nonlinear section strains to the element formulation. In the work by \cite{battini2002corotational-thesis} this resulted in a dedicated finite element implementation to incorporate the effect. \cite{lecorvec2012nonlinear} achieved a more modular implementation by averaging out the Wagner effect and incorporating it into a geometric transformation. However, this approach is less accurate, and remains coupled to the element implementation.
    In the present approach the effect is encapsulated in the section resultant model and is compatible with any finite element formulation in the family presented in \cref{sec:genwarp} without loss of accuracy as demonstrated in Example~\ref{sec:frame-0015} and Example~\ref{sec:frame-1035}.
\end{itemize}

The enhanced strain framework of \cref{sec:enhan-section} enriches the constitutive model without introducing auxiliary fields beyond those considered by the underlying beam theory, and is likewise encapsulated within the cross section resultant model. 
This has the limitation that it is only capable of applying corrections which vary along \(\reflenx\) in proportion with the strain measures \(\bm{e}\). 
However, this is the quality that allows the method to enhance response without introducing additional auxiliary fields. 

\iffalse
\paragraph{Remarks}
\begin{itemize}
\item Condition (i) ensures stability of the method for the linearized problem, and precludes interpolations for the enhanced strain field which would result in rank-deficient methods.
    \item 
Given an interpolation space $\tilde{\mathscr{E}}^h$ of enhanced strain fields, the orthogonality Condition~(ii) may produce an overly constrained space $\StressSpace^h$ which would preclude convergence of the method. 
We prevent such a situation by requiring Condition (iii).
\item For the linear theory, Condition~(iii) is closely related to satisfaction of the classical patch test (see \cite{taylor1986patch}).
\end{itemize}
\fi
